\chapter*{Résumé}
\addcontentsline{toc}{chapter}{Résumé} 
\vspace{-0.8cm}

Ce mémoire présente le travail réalisé lors de mon projet de fin d’études chez \textbf{NTT DATA}, dans le domaine de l’intelligence artificielle, dans le cadre du master \textbf{Systèmes Distribués et Intelligence Artificielle} à l’\textbf{ENSET de Mohammedia}. L’objectif principal était de concevoir un agent IA d’analyse des risques et de recommandation d’actions, fondé sur l’exploitation des données issues de \textbf{Jira Data Center} (projets, tickets, historiques). Le projet visait à automatiser la détection des anomalies, à prédire les retards et à proposer des actions correctives proactives. Plusieurs jeux de données ont été extraits, nettoyés et structurés. Des algorithmes de classification et des modèles de traitement du langage naturel ont été testés, intégrant divers critères (sévérité des tickets, temps de résolution, interactions entre équipes). Une interface de visualisation a également été développée pour simuler et ajuster les recommandations en temps réel. Ce rapport retrace les principales étapes du projet, de la préparation des données à la mise en œuvre des solutions, en soulignant les choix techniques et les résultats obtenus.
\vspace{2em}
\hline
\begin{flushright}
	\textbf{\textit{Mots-clés :}}NTT DATA, Agent IA, Jira Data Center, analyse des risques, recommandations, intelligence artificielle, NLP.
\end{flushright}
\hline